\section{Experimentelles Setup und Methodik}

\subsection{Simulationsumgebung}

Um die Fehler- und Manipulationsszenarien des Transport Layer Security (TLS) praktisch zu analysieren, wurde eine kontrollierte, lokale Loopback-Umgebung (127.0.0.1) auf einem Windows-Endgerät implementiert. Als kryptographische Grundlage kam OpenSSL (Version 3.6.1) sowohl auf Client- als auch auf Server-Seite zum Einsatz~\cite{openssl}. Aufgezeichnet wurde der Datenverkehr parallel auf dem Loopback-Interface mittels der Netzwerk-Analysesoftware Wireshark, um den exakten Fluss der Handshake-Pakete zu dokumentieren. Außerdem wurden aktuelle Konfigurationsrichtlinien für sichere Web-Umgebungen herangezogen~\cite{mdndocs}.


\subsection{TLS-Fehlersimulation und Log-Analyse}
Die Aufgabenstellung erfordert die Simulation eines Downgrade- bzw.\ Handshake-Fehlerfalls sowie die Analyse der Stellen, an denen Schutzmechanismen greifen oder geschwächt werden.


\subsection{Szenario 1: Nicht-erfolgreicher Downgrade (Handshake-Fehlerfall)}
Im ersten Experiment wurde untersucht, wie ein moderner Client reagiert, wenn ein Server ein veraltetes, unsicheres Protokoll erzwingt. Dies simuliert das Verhalten bei einer serverseitigen Fehlkonfiguration oder einem Manipulationsversuch im Netzpfad~\cite{sentinelone}.

\subsection*{Vorbereitung der Zertifikate}
Zunächst wurde durch die Nutzung von OpenSSL ein selbstsigniertes X.509-Zertifikat mit einem 2048-Bit RSA-Schlüsselpaar generiert. Der Parameter \texttt{-nodes} (No-DES) sorgt dafür, dass der private Schlüssel unverschlüsselt bleibt, um einen reibungslosen Serverstart ohne Passwortabfrage zu ermöglichen:

\begin{lstlisting}[style=consolestyle]
openssl req -x509 -newkey rsa:2048 -keyout key.pem -out cert.pem \
  -days 365 -nodes
\end{lstlisting}

\subsection*{Verbindungsaufbau}
Der simulierte Server wurde so gestartet, dass er ausschließlich das veraltete Protokoll TLS~1.0 anbietet:

\begin{lstlisting}[style=consolestyle]
openssl s_server -accept 4433 -cert cert.pem -key key.pem -tls1
\end{lstlisting}

Anschließend versuchte der reguläre Client, eine sichere Verbindung zu diesem Port aufzubauen:

\begin{lstlisting}[style=consolestyle]
openssl s_client -connect localhost:4433
\end{lstlisting}

\subsection*{Analyse der Handshake-Schritte}

Der Verbindungsaufbau wird durch den Client mit dem \textit{Client Hello} gestartet, das standardmäßig die Unterstützung moderner Protokolle (wie TLS~1.3) signalisiert. Der Server antwortet stur mit einem \textit{Server Hello} und versucht, TLS~1.0 aufzuzwingen~\cite{ssldragon}. An dieser Stelle greift der clientseitige Sicherheitsmechanismus: Moderne OpenSSL-Bibliotheken blockieren standardmäßig die Kommunikation über Protokolle unterhalb von TLS~1.2~\cite{openssl}. Der Handshake bricht unmittelbar ab. Das Terminal des Clients dokumentiert diesen fatalen Protokollkonflikt wie folgt:

\begin{lstlisting}
CONNECTED(000001B8)
FC270000:error:0A00042E:SSL routines:ssl3_read_bytes:
tlsv1 alert protocol version:../openssl-3.6.1/ssl/record/
rec_layer_s3.c:918:SSL alert number 70
---
no peer certificate available
\end{lstlisting}

In der parallelen Wireshark-Aufzeichnung zeigt sich dieser Abbruch durch ein rot markiertes Paket mit der Kennzeichnung \texttt{Alert (Level: Fatal, Description: Protocol Version (70))}. Der Fehlercode~70 (\textit{protocol\_version}) belegt, dass der Client die erzwungene Abwärtskompatibilität des Servers zum Schutz des Anwenders blockiert hat~\cite{alertcode70}. In echten Systemen verhindert dieser Mechanismus, dass Verbindungen unbemerkt auf unsichere Standards zurückfallen.

\subsection{Szenario 2: Erfolgreicher Downgrade (Schwächung der Mechanismen)}
Das zweite Experiment simuliert den Fall, in dem die clientseitigen Schutzmechanismen bewusst deaktiviert werden -- beispielsweise, um die Kommunikation mit ungenügend gehärteten Altsystemen (Legacy-Systemen) im Firmennetzwerk zu erzwingen~\cite{sentinelone}.

\subsection*{Verbindungsaufbau unter reduzierten Sicherheitsbarrieren}
Nun werden der Server und der Client durch das explizite Herabsetzen des Sicherheitslevels auf die Stufe Null (\texttt{@SECLEVEL=0}) jeweils gezwungen, ihre Schutzschilde zu deaktivieren~\cite{openssl}:

\begin{lstlisting}[style=consolestyle]
# Server:
openssl s_server -accept 4433 -cert cert.pem -key key.pem -tls1
  -cipher "DEFAULT@SECLEVEL=0"

# Client:
openssl s_client -connect localhost:4433 -tls1
  -cipher "DEFAULT@SECLEVEL=0"
\end{lstlisting}

\subsection*{Analyse der Handshake-Schritte}
Durch die Modifikation signalisiert der Client bereits im \textit{Client Hello}, dass er im Notfall auch das unsichere TLS~1.0 akzeptiert~\cite{ssldragon}. Der Server antwortet im \textit{Server Hello} mit der Bestätigung der veralteten Version. Da der Client das nicht mehr blockiert, wird der Handshake ohne Fehlermeldung abgeschlossen. Die Verbindung wechselt in den Zustand \texttt{CONNECTED}. Das Terminal des Clients liefert folgendes Protokoll-Log:

\begin{lstlisting}
CONNECTED(000001BC)
Can't use SSL_get_servername
depth=0 C=AU, ST=Some-State
verify error:num=18:self-signed certificate
verify return:1
---
Certificate chain
 0 s:C=AU, ST=Some-State
   i:C=AU, ST=Some-State
   a:PKEY: RSA, 2048 (bit); sigalg: sha256WithRSAEncryption
---
New, TLSv1.0, Cipher is ECDHE-RSA-AES256-SHA
Protocol  : TLSv1
Server public key is 2048 bit
Secure Renegotiation IS NOT supported
\end{lstlisting}

\subsection*{Bewertung der Schwachstelle}
An dieser Stelle versagen die Sicherheitsmechanismen der Transportschicht im realen Betrieb~\cite{sentinelone}. Obwohl für den Anwender eine etablierte Verschlüsselung signalisiert wird (\texttt{CONNECTED}), ist die Vertraulichkeit und Integrität der Daten nach heutigem Standard stark gefährdet. Da TLS~1.0 (im Log als \texttt{TLSv1} ausgewiesen) strukturelle Schwachstellen besitzt und keine modernen Authentifizierungsverfahren (wie AEAD) erzwingt, können Angreifer im Netzpfad (Man-in-the-Middle) die Sitzung mit modernen Rechenkapazitäten passiv abhören oder manipulieren~\cite{sentinelone, mdndocs}.

\subsection{Analytische Gegenüberstellung der Simulationsergebnisse}
Zur übersichtlichen Strukturierung der Ergebnisse sind die Auswirkungen der beiden Konfigurationen in Tabelle~\ref{tab:vergleich} zusammenfassend gegenübergestellt.

\begin{table}[ht]
\caption{Vergleichsmatrix der simulierten TLS-Fehlerfälle}
\label{tab:vergleich}
\centering
\footnotesize
\renewcommand{\arraystretch}{1.3}
\resizebox{\columnwidth}{!}{
\begin{tabular}{lll}
\toprule
\textbf{Kriterium}
  & \textbf{Szenario 1: Mismatch}
  & \textbf{Szenario 2: Downgrade} \\
\midrule
Server-Vorgabe
  & Erzwingung TLS~1.0
  & Erzwingung TLS~1.0 \\
Client-Konfiguration
  & Standard (Level $\geq$ 1)
  & Deaktiviert (Level~0) \\
Handshake-Verlauf
  & Abbruch (Fatal Alert)
  & Erfolgreich (\texttt{CONNECT}) \\
Log-Indikator
  & \texttt{SSL alert number 70}
  & \texttt{Protocol: TLSv1} \\
Sicherheitszustand
  & Sicher (Leitung blockiert)
  & Insecure (vulnerabel) \\
\bottomrule
\end{tabular}
}
\end{table}